% !TEX root = ../main.tex
\section{Hallucination as Unwarranted Assertion}
\label{sec:hallucination}

\subsection{The definitions}

\begin{definition}[Hallucination and error]
\label{def:hallucination}
\begin{align}
  \Hall(\phi \mid c) &\;=_{\mathrm{def}}\; \Out_M \phi \wedge \neg\,\Grd_M(\phi \mid c),
    \label{eq:hallucination}\\
  \Err(\phi \mid c)  &\;=_{\mathrm{def}}\; \Out_M \phi \wedge \Grd_M(\phi \mid c) \wedge \neg \phi.
    \label{eq:error}
\end{align}
\end{definition}

Observe what \eqref{eq:hallucination} omits: any truth-value condition.
Hallucination is unwarranted assertion---assertion unsupported by evidence
applicable to the context of use---and that is all. Truth and falsity classify
errors and successes; they do not classify hallucinations. The argument for this
regimentation runs through the off-diagonal cells and the applicability
distinction.

Before making it, I should place the account. The distinction between consistency
with a source and correspondence with the world is not new to the technical
literature, where it appears as \emph{faithfulness} against \emph{factuality}
\citep{maynez2020,ji2023}. My grounded/true grid is recognisably that
distinction, and I claim no priority for it. What I claim is threefold: that the
normative asymmetry runs one way---only ungroundedness constitutes
hallucination---that the defect in the ungrounded-but-true case admits a precise
modal analysis, and that the distinction follows from a theory of warranted
assertion rather than from engineering convenience. The account is also close
kin to the argument that these systems produce assertion without regard for
truth in Frankfurt's sense \citep{hicks2024}; \eqref{eq:hallucination} may fairly
be read as a formalisation of that claim, with the difference that it is
indexed, measurable, and yields a metric.

\subsection{Faithful transmission of a bad source is error, not hallucination}
\label{sec:error}

Suppose $M$'s corpus contains, pervasively and consistently, a claim that a
treaty was signed in 1748 when it was signed in 1749---a common error propagated
across reference works. $M$ asserts the 1748 date, and attribution analysis
confirms the assertion is driven by exactly those contents. On the
truth-conditional definition this is a hallucination: false output. But that
classification obliterates a distinction we need. $M$ has done everything a
testimonial transmitter can do: it reported what its evidence says, through a
mechanism sensitive to that evidence. The failure lies in the evidence, not in
the assertion mechanism.

The distinction is not pedantic; it is intervention-relevant, and that is the
form the argument should take for anyone who has to allocate engineering effort.
Errors are corrected by fixing the evidential base: better data, retrieval from
authoritative sources, corpus curation. Hallucinations are corrected by fixing
the assertion mechanism: making generation conditional on evidential support, and
abstaining when support is absent. A metric that pools cells (iii) and (iv) as
``hallucination rate'' tells an engineer nothing about which intervention is
needed, and in my experience the two are handled by different teams on different
budgets. Normatively, too: we do not say that a scrupulous newspaper
hallucinated when it accurately reported a source that turned out to be wrong.
We say it erred, or was misled. The same discipline should govern machine
assertion.

\subsection{Ungrounded truth is a safety failure, not a Gettier case}
\label{sec:veridical}

Now the converse. $M$ is asked for the population of a small town absent from its
corpus. The generation mechanism, running on formulation-level
statistics---towns with names of this shape, in sentences of this genre, get
populations of this magnitude---produces ``8,200.'' By luck, that is correct. On
the truth-conditional definition nothing has gone wrong. But everything
distinctive of hallucination has occurred. The mechanism was disconnected from
evidence about the town, and the truth of the output is an accident.

The precise nature of the defect is modal. Had the world been
different---had the population been 3,400---the very same process would have
produced the very same output. An assertion with that profile does not track the
fact it concerns; it is unsafe, in the sense of \citet{sosa1999}, and the
correct label for the case is \emph{veritic luck} \citep{pritchard2005}. An
assertion so produced is defective as an assertion regardless of its truth value.
I will call it a \emph{veridical hallucination}, and the terminology should mark
that the defect in cell (ii) and cell (iv) is one and the same defect, differing
only in worldly luck.

It is tempting to call this a Gettier structure, and I want to resist the
temptation precisely because something Gettier-like is going on nearby.
\citet{gettier1963} showed that \emph{justified} true belief can fall short of
knowledge when luck severs the connection between the justification and the
truth. In cell (ii) the justification is not severed but absent: there is no
grounding at all, only truth. That is the pre-Gettier case---true belief without
justification---and the lesson it teaches was already learned when justification
was added to the analysis.

The genuine Gettier structure appears one step downstream, in the recipient. A
user receives a fluent, well-formed, contextually apt assertion. That fluency
functions, reasonably enough, as apparent justification: it is the signal humans
have always used to calibrate trust in testimony. The content is true. So the
user forms a belief that is true and, by her lights, justified---and it is not
knowledge, because what justified it stands in no connection to what made it
true. Cell (ii) is a Gettier-case generator, and this is why it matters
practically rather than taxonomically. A user who receives a cell-(ii) output and
verifies it acquires no reason to trust the system's next output, though
truth-conditional evaluation will have scored the system a success. Lucky truths
are not repeatable performances. Grounding is.

I note that condition (2) of Definition~\ref{def:grounding} and the analysis just
given are the same condition applied twice. Both demand that the assertion event
be counterfactually sensitive to something---to the evidence in the first case,
to the fact in the second. The account is accordingly more unified than the
section headings suggest: hallucination is assertion that fails a sensitivity
condition on its evidential base, and veridical hallucination is that failure
masked by luck.

\subsection{Misapplied grounding: the third failure class}
\label{sec:misapplied}

The context index makes visible a failure that neither cell (iii) nor cell (iv)
captures, and it is the one that dominates regulated deployment.

A model asserts that a compound is approved for an indication. The assertion is
driven by documents in $E$ that say exactly that, and it is counterfactually
sensitive to them: remove them and the assertion goes away. But the approval was
withdrawn in 2024, and the decision at hand is being taken now. Or the documents
describe one jurisdiction's regime and the decision falls under another's. Or the
evidence concerns an adult population and the intended use is paediatric.

On the unindexed definition these are all cell (iii): grounded, false, therefore
error, therefore charged to a defective evidential base. That verdict is wrong
twice over. The evidential base was not defective---the 2019 documents were
accurate in 2019, and the foreign regime is correctly described. And the remedy
implied is wrong: curating the corpus will not help, because nothing in it is
incorrect. What failed is applicability, and the remedy is
resolution---establishing $c$ before answering, checking currency at $\tau$,
checking jurisdiction $j$ and intended use $u$, and abstaining or qualifying when
the applicable evidence has not been identified.

On Definition~\ref{def:grounding} these cases are ungrounded relative to $c$, and
so fall on the hallucination side. I take that to be the right result: evidence
that does not bear on the decision at hand provides no support for a claim
advanced for that decision, however authentic it was in its own context. The
practical upshot is a three-way partition of defective assertion, with three
distinct remedies:

\begin{center}
\begin{tabular}{@{}lll@{}}
\toprule
Failure & Diagnosis & Remedy \\
\midrule
Error & applicable evidence, but wrong & corpus and source curation \\
Misapplied grounding & evidence inapplicable at $c$ & currency and scope resolution \\
Fabrication or veridical luck & no evidential connection & abstention, grounding gates \\
\bottomrule
\end{tabular}
\end{center}

A hallucination metric that cannot separate these three tells you which of three
budgets to spend nothing about.

\subsection{Accuracy places no bound on grounding}
\label{sec:independence}

The off-diagonal arguments jointly indict the dominant evaluation methodology.
Truth-matching benchmarks misclassify cell (iii) as hallucination, penalising
faithful transmission and thereby rewarding systems that override their evidence
when the evidence is unfashionable; and they misclassify cell (ii) as success,
rewarding lucky fabrication and thereby favouring confident guessing over
abstention. The misclassifications are not symmetric noise. Both push
optimisation pressure the same way---toward assertion mechanisms decoupled from
evidential support and tuned to the benchmark's answer distribution.

The relationship can be stated exactly. Let
\[
  a = \Pr\bigl(\phi \mid \Out_M \phi\bigr), \qquad
  g = \Pr\bigl(\Grd_M(\phi \mid c) \mid \Out_M \phi\bigr)
\]
be the accuracy and the grounding rate of a system's assertions.

\begin{proposition}
\label{prop:independence}
Every pair $(a, g) \in [0,1]^2$ is realisable.
\end{proposition}

\begin{proof}
Let the cells of Table~\ref{tab:taxonomy} carry masses
$m_{\mathrm{GT}}, m_{\mathrm{GF}}, m_{\mathrm{UT}}, m_{\mathrm{UF}}$ summing to
$1$, so that $a = m_{\mathrm{GT}} + m_{\mathrm{UT}}$ and
$g = m_{\mathrm{GT}} + m_{\mathrm{GF}}$. Given a target $(a,g)$, choose
\[
  m_{\mathrm{GT}} \in \bigl[\max(0,\, a + g - 1),\ \min(a, g)\bigr],
\]
an interval that is non-empty for every $(a,g)$ in the unit square. Set
$m_{\mathrm{GF}} = g - m_{\mathrm{GT}}$,
$m_{\mathrm{UT}} = a - m_{\mathrm{GT}}$, and
$m_{\mathrm{UF}} = 1 - a - g + m_{\mathrm{GT}}$. Each is non-negative by the
choice of interval, and the four sum to $1$.
\end{proof}

The construction requires that grounding not entail truth, which cell (iii)
grants: an assertion may be grounded in applicable evidence and false. Given
that, accuracy and grounding vary independently.

\begin{corollary}
\label{cor:no-bound}
No truth-matching metric bounds the hallucination rate $1 - g$. A benchmark
reporting accuracy alone is consistent with any hallucination rate whatever.
\end{corollary}

The result is elementary, and I state it because the elementary version is what
the field's practice currently violates. A second consequence is worth drawing:
since $a$ and $g$ are free to move independently, accuracy can improve while the
hallucination rate worsens. That is a plausible reading of what optimising
against answer keys does, and it is not a hypothetical concern for anyone
choosing between checkpoints on a leaderboard.

\subsection{The measurement target, and where it can be measured}
\label{sec:measurement}

The correct target is the conditional probability that an assertion is
ungrounded given that it is made:
\begin{equation}
  \text{hallucination rate} \;=_{\mathrm{def}}\;
  \Pr\bigl(\neg\,\Grd_M(\phi \mid c) \;\big|\; \Out_M \phi\bigr).
  \label{eq:rate}
\end{equation}
This is estimated by attribution audit rather than answer-key comparison, and
the framework for doing so already exists in the work on attributing generated
statements to identified sources \citep{rashkin2023}.

One caveat is essential, and it is the point at which this paper's argument
begins to bear on system design rather than only on measurement. The two halves
of $E$ are not equally tractable. For retrieval-provided context, condition (2)
of Definition~\ref{def:grounding} is directly testable: ablate or contradict the
retrieved evidence and observe whether the assertion survives. For parametric
encodings it is not. ``Had $E$ not supported $\phi$'' ranges over counterfactual
training corpora, and evaluating it properly means retraining.
Training-data attribution methods approximate the counterfactual, but at
frontier scale they remain an open problem, and the difference between the two
cases is one of kind rather than degree.

So \eqref{eq:rate} is well defined for any system and estimable only for systems
that make their evidential base explicit at inference time. An unaugmented
model's hallucination rate is a perfectly good quantity that nobody can currently
measure. This is not a defect in the definition; it is an argument about
architecture, and it is the strongest one this paper has. If the rate at which a
system asserts beyond its applicable evidence is the quantity that matters for
consequential use, then systems must be built so that the quantity can be
observed---with an identified, versioned, context-resolved evidential base at the
point of assertion, not merely a training corpus somewhere behind it.

\subsection{Why the two contributions are one}
\label{sec:unity}

The two halves of this paper are not independent. Why do these systems
hallucinate in the defined sense? Because $\Out_M$ is hyperintensional in the way
Section~\ref{sec:re-failure} diagnosed. Assertion probability is computed over
formulations, and a formulation can be made probable by surface statistics with
no route through evidential support for its content. There is also a statistical floor here that no amount of architectural
care removes: \citet{kalai2024} show that a calibrated model must hallucinate at
some positive rate on facts appearing rarely in its evidence, which is a formal
reason to expect the ungrounded cells to be non-empty rather than merely badly
managed. Cells (ii) and (iv) are the
behavioural face of RE-failure: cases where the string-level machinery runs free
of the content-level relation $\Grd_M$ tracks. Conversely, a system whose
assertions were closed under the right relations to its applicable evidence
would, to that extent, approximate normality. The formal diagnosis and the
redefinition of hallucination are two descriptions of one architecture.
